\chapter{Conclusiones}

% NOTA: capítulo a completar al finalizar el desarrollo. Debe recoger una valoración del trabajo
% realizado, el grado de cumplimiento de los objetivos planteados en la introducción y una
% reflexión personal, además de las líneas de mejora y trabajo futuro.

\section{Experiencia personal}

% Reflexión sobre el proceso de desarrollo: aprendizajes técnicos (stack web moderno, visualización
% 3D con WebGL, contenedores), dificultades encontradas y su resolución, y valoración global del
% cumplimiento de los objetivos.

\section{Propuesta de mejoras}

Entre las principales líneas de trabajo futuro se contemplan:

\begin{itemize}
	\item Optimización del rendimiento para grafos de gran tamaño (más de 1000 nodos) mediante
		técnicas como el nivel de detalle (LOD) o el uso de \textit{web workers} para los cálculos
		más pesados.
	\item Incorporación de análisis de red más avanzados, como la detección de comunidades o el
		cálculo de caminos mínimos entre regiones.
	\item Soporte para más atlas y formatos de datos públicos.
	\item Autenticación real y persistencia de configuraciones de usuario.
\end{itemize}
